% Color RGBA: 0c4a6eff (RGB: #0c4a6e, fully opaque)
\documentclass[12pt]{article}
\usepackage{amsmath}
\usepackage{amssymb}
\usepackage{xcolor}

\definecolor{customcolor}{HTML}{0c4a6e}
\pagecolor[HTML]{FFFFFF}
\color{customcolor}

\begin{document}
\thispagestyle{empty}
\setlength{\parindent}{0pt}
\centering

{\Large\bfseries ¿Por qué aparecen los términos ``extra'' en la aceleración polar?}

\vspace{0.5em}

$\frac{\Delta\hat{e}_r}{\Delta t} \approx \frac{\Delta\varphi}{\Delta t}\cdot\hat{e}_\varphi \rightarrow \frac{d\hat{e}_r}{dt} = \dot{\varphi}\cdot\hat{e}_\varphi$ del mismo modo: $\frac{d\hat{e}_\varphi}{dt} = -\dot{\varphi}\cdot\hat{e}_r$

\vspace{0.3em}

Al derivar $\dot{r} = \dot{r}\cdot\hat{e}_r + r\dot{\varphi}\cdot\hat{e}_\varphi$ → el término $-r\dot{\varphi}^2$ viene de: $r\dot{\varphi}\cdot\left(\frac{d\hat{e}_\varphi}{dt}\right) = r\dot{\varphi}\cdot(-\dot{\varphi}\cdot\hat{e}_r)$

\end{document}
